% Session 7 -- Refusals and Faults: Procedure TPS-P-08, Checklist TPS-C-3.
% Law 1 sources: spec/tps_ritual.json (every refusal and its exit code:
% TPS-001/002 exit 1, TPS-003/004 ask again, TPS-005 exit 5, TPS-006 exit
% 6, TPS-007 exit 7, TPS-008 exit 8); os/milton/tpsrep_program.asm (the
% table is opened after the parameter record and before the alignment, so
% TPS-006/007 come before any ritual; TPS-006 names the first required field
% missing, or SIGNATURE when the file is not a dBASE III table; TPS-008 after
% the attestation, when TPSREPT.TXT cannot be created); spec/dos_messages.json
% and os/milton/int21.c (the device error messages and Abort, Retry, Fail?,
% as the InitechDOS Reference Appendix A documents them); COMMAND.COM's Bad
% command or file name (MSG-DOS-0002).  The Help Desk is extension 2504
% (README.md, PUB-REF-0001 D-12).  The program issues no printer I/O: this
% session documents no printer fault (research brief Sec 4 is unapproved).
% Screens tps007, tps006, tpsroot: the journey route.
% GATES: test-tps-report test-flair-tps-journey test-pubs-screens test-pubs-tps test-dosmsg
\Gates{test-tps-report, test-flair-tps-journey, test-pubs-screens, test-pubs-tps, test-dosmsg}
\chapter{Refusals and Faults}

\SideHead{Introduction}The program refuses a step that is not carried out as
the procedure requires.  Some refusals ask again; the others end the run
with a return code.  This session lists every refusal, where it occurs, and
the action it requires, and gives the procedure for starting again and the
checklist for escalation.

\SideHead{Objectives}Upon completion of this session you will be able to:
\begin{enumerate}
\item Name the refusal that each step of the operating procedure can give.
\item Start a run again after a refusal that ends it.
\item Respond to an InitechDOS device error message.
\item Escalate as Checklist TPS-C-3 requires.
\end{enumerate}

\section{The Refusals}
\index{refusals}
The program checks its input in this order.  The appendix of messages
(Appendix G) gives each message in full.

\begin{center}\small\hspace*{-\imHang}%
\begin{tabular}{@{}p{30mm}p{40mm}p{24mm}p{\dimexpr\imFullWidth-94mm-6\tabcolsep\relax}@{}}
\rule{0pt}{10pt}\sffamily\bfseries Message & \sffamily\bfseries When & \sffamily\bfseries Result & \sffamily\bfseries Action\\\hline
TPS-001 & No parameter record & Ends, code 1 & \proc{01} step 6; else supervisor\\
TPS-002 & A line of the record missing & Ends, code 1 & Supervisor (OP-7.3)\\
TPS-007 & No table named, or not found & Ends, code 7 & Type the command again\\
TPS-006 & The table is not a TPS table & Ends, code 6 & Supervisor\\
\texttt{FORM NOT ALIGNED} & Wrong number of presses & Asks again & \proc{03} step 1\\
TPS-003 & Report numbers differ & Asks again & \proc{04} step 1\\
TPS-004 & Not the memo number & Asks again & \proc{05} step 2\\
TPS-005 & Cover sheet answer N & Ends, code 5 & \proc{06}; run again\\
TPS-008 & Report file not created & Ends, code 8 & Supervisor\\\hline
\end{tabular}\end{center}

\section{Refusals Before the Alignment}
After the parameter record, the program opens the table named on the command
line.  If none is named, or it cannot be opened, the program displays its
usage:

\OpsShot[fromtext={A:\Bsl APPS>TPSREP}, caption={The command without a table
name.}]{tps007}

If the file named is not a TPS table, the program names the first of the six
fields it does not find, or \texttt{SIGNATURE} when the file is not a dBASE
III table at all:

\OpsShot[fromtext={A:\Bsl APPS>TPSREP}, caption={A file that is not a dBASE
III table, named as the table.}]{tps006}

\section{Starting Again}
\begin{Procedure}{TPS-P-08}{To Start a Run Again After a Refusal or a Fault}
\ProcField{Purpose}{To end a refused or interrupted run in order and start it
  again from the beginning of the operating procedure.}
\ProcField{Frequency}{Each time a run ends with a return code other than 0,
  or a device error message is displayed.}
\ProcField{Performed by}{The operator; the supervisor where Checklist TPS-C-3
  requires.}
\ProcField{Prerequisites}{The prompt \texttt{A:\bs APPS>} is displayed, or a
  device error message.}
\ProcField{Materials}{\form{5}; Checklist TPS-C-3.}
\begin{ProcSteps}
\item If a device error message and \texttt{Abort, Retry, Fail?} are
  displayed, correct the cause named in the message and press \type{R}.  If
  the message is displayed again, press \type{A}, and go to step 5.
\item Read the end line \texttt{TPS REPORT GENERATOR TERMINATED.} and note the
  return code.
\item Enter the message number, the return code, the date and the time on
  \form{5}, line 2.
\item Look up the message in the table above.  If the action is to see your
  supervisor, go to step 5.  Otherwise take the action it gives.
\item Complete Checklist TPS-C-3.
\item When the cause has been corrected, start again at \proc{02}, step 1.
  Each step of the operating procedure is carried out again in full, the
  alignment included.
\end{ProcSteps}
\ProcField{Verification}{The new run ends with \texttt{RETURN CODE 0}.}
\ProcField{Records retained}{\form{5}, line 2; Checklist TPS-C-3, if
  completed; the forms of the new run.}
\ProcNote{A refusal that asks again (\texttt{FORM NOT ALIGNED}, TPS-003,
  TPS-004) does not end the run.  Continue at the step the table gives.}
\end{Procedure}

\section{InitechDOS Device Error Messages}
\index{device error messages}
When InitechDOS cannot read or write the disk, it displays a device error
message, followed by \texttt{Abort, Retry, Fail?}.  Press one letter; you need
not press \key{Enter}.  \type{R} tries the operation again; \type{A} ends the
program and returns to the prompt; \type{F} tells the program that the
operation failed.  The device error messages are listed in Appendix~G and,
with the rest of the InitechDOS messages, in the \doctitle{InitechDOS
Reference}, Appendix A.

When the command itself is not found, InitechDOS displays \texttt{Bad command
or file name}.  At the prompt \texttt{A:\bs APPS>} this means that a program
file is missing: see your supervisor.  At any other prompt, see
\proc{01}, step 6.

\section{Escalation}
\index{escalation}\index{supervisor}
\begin{Checklist}{TPS-C-3}{Escalation}
\CheckItem{The message, the return code, the date and the time are on
  \form{5}, line 2.}
\CheckItem{The run was refused at \texttt{A:\bs APPS>}, not at another prompt.}
\CheckItem{The program has not been run again since the refusal, except as
  the table of refusals directs.}
\CheckItem{Your supervisor has been told, and has initialled \form{5},
  line 2.}
\CheckItem{If your supervisor directs, the Help Desk (extension 2504) has
  been called, and the call is noted on \form{5}, line 2.}
\end{Checklist}

\section{Lab Project 6. Refusals}
\begin{enumerate}
\item At the prompt \texttt{A:\bs APPS>}, type \type{TPSREP} and press
  \key{Enter}.  Note the message and the return code:\par
  \rule{40mm}{0.4pt}
\item Type \type{TPSREP TPSPARM.DAT} and press \key{Enter}.  Note the field
  named: \rule{30mm}{0.4pt}
\item Carry out \proc{08}, steps 2 to 4, for each refusal, on a practice copy
  of \form{5}.
\end{enumerate}

\section{Study Questions}
\begin{enumerate}
\item Which refusals ask again, and which end the run?
\item What return code follows the answer \type{N} to the cover sheet
  prompt?
\item After a refusal that ends the run, from which procedure do you start
  again?
\item What do you press when \texttt{Abort, Retry, Fail?} is displayed after
  you have corrected the cause?
\item What is the extension of the Help Desk?
\end{enumerate}
